\documentclass[11pt,a4paper]{article}
\usepackage[margin=1in]{geometry}
\usepackage{amsmath}
\usepackage[numbers,sort&compress]{natbib}
\usepackage{hyperref}
\usepackage{booktabs}
\usepackage[french]{babel}

\title{Les contributions techniques de W.B.\ Hutchinson \\
       \`a la m\'etallurgie physique}
\author{R\'esum\'e technique}
\date{Mars 2026}

\begin{document}
\maketitle

\begin{abstract}
W.B.\ (Bevis) Hutchinson, ayant effectu\'e l'essentiel de sa carri\`ere
\`a l'Institut su\'edois de recherche sur les m\'etaux (SIMR, devenu
Swerea KIMAB puis Swerim) \`a Stockholm, a apport\'e des contributions
fondamentales \`a la compr\'ehension de la texture cristallographique,
des m\'ecanismes de recristallisation et des relations
structure--propri\'et\'es dans les aciers et les alliages d'aluminium.
Avec plus de 180 publications et plus de 7\,000 citations, ses travaux
couvrent les aciers doux pour t\^oles, les aciers sans interstitiels
(IF), les aciers \'electriques \`a grains orient\'es, les aciers
microalli\'es de construction et les produits en t\^ole d'aluminium.  Ce
document r\'esume les principaux th\`emes techniques de ses recherches.
\end{abstract}

%% ===================================================================
\section{D\'eveloppement de la texture dans les aciers doux pour t\^oles}
\label{sec:bcc-texture}

La revue de Hutchinson parue en 1984 dans \emph{International Materials
Reviews}, \og Development and control of annealing textures in low-carbon
steels \fg{} \cite{Hutchinson1984IMR}, est l'un de ses travaux les plus
influents.  Elle traite de la formation de la texture de fibre~$\gamma$
($\{111\}\langle uvw\rangle$) qui se d\'eveloppe lors du recuit de
recristallisation des t\^oles d'acier ferritique lamin\'ees \`a froid.

\subsection{La fibre~$\gamma$ et l'aptitude \`a l'emboutissage}

Le coefficient de Lankford $\bar{r}$ (rapport d'anisotropie plastique
normal) gouverne l'aptitude \`a l'emboutissage profond des t\^oles
d'acier.  Des valeurs \'elev\'ees de $\bar{r}$ sont corr\'el\'ees
\`a une fibre~$\gamma$ intense, c'est-\`a-dire une pr\'edominance de
grains dont les plans $\{111\}$ sont parall\`eles \`a la surface de la
t\^ole.  La composante de texture concurrente est
$\{100\}\langle 0vw\rangle$ (fibre~$\alpha$), qui d\'egrade
l'emboutissabilit\'e.  Hutchinson a montr\'e que le rapport des
intensit\'es $\{111\}$/$\{100\}$ dans la texture de recristallisation
est le facteur d\'eterminant pour la valeur de $\bar{r}$, et que ce
rapport peut \^etre manipul\'e par la chimie de l'alliage (carbone,
azote, mangan\`ese et \'el\'ements de microalliage), la microstructure
de la bande \`a chaud, le taux de r\'eduction au laminage \`a froid et
les conditions de recuit (recuit en caisse \emph{vs.}\ recuit continu).

\subsection{Nucl\'eation \`a partir des grains de fibre~$\gamma$}

Dans l'acier doux d\'eform\'e, les grains
$\{111\}\langle 112\rangle$ d\'eveloppent des bandes de cisaillement
denses lors du laminage \`a froid.  Hutchinson a identifi\'e ces bandes
de cisaillement comme des sites de nucl\'eation pr\'ef\'erentiels pour
les grains recristallis\'es qui h\'eritent de l'orientation $\{111\}$.
Cette \og nucl\'eation orient\'ee \fg{} \`a partir des grains de
fibre~$\gamma$ cisaill\'es est le m\'ecanisme principal par lequel une
texture de recristallisation \`a fibre~$\gamma$ intense s'\'etablit,
conduisant \`a des valeurs \'elev\'ees de $\bar{r}$.

\subsection{Aciers sans interstitiels (IF)}

Le d\'eveloppement des aciers IF (stabilis\'es au Ti et/ou au Nb pour
pi\'eger tout le C et le N interstitiels) a cr\'e\'e un v\'ehicule
id\'eal pour l'optimisation de la texture de fibre~$\gamma$.  En
l'absence d'atmosph\`eres de Cottrell et de pr\'ecipit\'es fins de
carbures pour interf\'erer avec la migration des joints de grains, les
aciers IF d\'eveloppent des textures de recristallisation $\{111\}$
extr\^emement fortes et atteignent des valeurs de $\bar{r}$
sup\'erieures \`a~2.  Le cadre d'analyse de texture de Hutchinson
sous-tend la conception des proc\'ed\'es industriels pour ces aciers.

%% ===================================================================
\section{M\'ecanismes de recristallisation et sous-structure}
\label{sec:rex}

Un th\`eme r\'ecurrent dans les travaux de Hutchinson est le lien entre
la sous-structure de d\'eformation et la nucl\'eation de la
recristallisation.

\subsection{Sous-structures de d\'eformation}

Dans une revue publi\'ee dans \emph{Materials Science Forum}
\cite{Hutchinson2007MSF}, Hutchinson a pass\'e en revue les
sous-structures de dislocations cr\'e\'ees par la d\'eformation
plastique des m\'etaux, en portant une attention particuli\`ere aux
sous-structures qui accommodent des d\'esorientations marqu\'ees.  Ces
\'el\'ements \`a forte d\'esorientation---bandes de d\'eformation,
bandes de transition, bandes de cisaillement et r\'egions proches des
joints de grains o\`u les rotations du r\'eseau s'accumulent---rev\^etent
une importance fondamentale car ils fournissent les joints de grains \`a
grand angle n\'ecessaires \`a la nucl\'eation de la recristallisation.

Diff\'erents modes de d\'eformation et diff\'erentes orientations
cristallines produisent diff\'erents types de sous-structure, de sorte
que le m\'ecanisme de nucl\'eation et la texture de recristallisation
r\'esultante d\'ependent du chemin de d\'eformation ant\'erieur.  Cette
compr\'ehension donne un levier pratique~: en contr\^olant les conditions
de d\'eformation (temp\'erature, d\'eformation, trajet de d\'eformation),
on peut influencer les m\'ecanismes de nucl\'eation actifs et ainsi
orienter la texture recristallis\'ee.

\subsection{Le d\'ebat nucl\'eation orient\'ee \emph{vs.}\ croissance s\'elective}

L'origine des textures de recristallisation a \'et\'e d\'ebattue
pendant des d\'ecennies.  Deux hypoth\`eses classiques sont~:
(i)~la nucl\'eation orient\'ee---de nouveaux grains nucl\'eent
pr\'ef\'erentiellement \`a partir de certaines orientations~; et
(ii)~la croissance s\'elective---les grains de certaines orientations
b\'en\'eficient d'un avantage de mobilit\'e et consomment la matrice
d\'eform\'ee plus rapidement.  Les travaux exp\'erimentaux de
Hutchinson, utilisant notamment le HVEM (microscopie \'electronique \`a
haute tension) pour observer la nucl\'eation in situ, ont apport\'e des
preuves en faveur de la nucl\'eation orient\'ee comme facteur dominant
dans de nombreux syst\`emes, bien qu'il ait reconnu que les deux
m\'ecanismes peuvent op\'erer et qu'un mod\`ele combin\'e est
n\'ecessaire pour une pr\'ediction quantitative.

%% ===================================================================
\section{Texture cube dans les m\'etaux CFC}
\label{sec:cube}

Dans les m\'etaux CFC (aluminium, cuivre, nickel), un laminage \`a froid
intense produit une texture de laminage de fibre~$\beta$, et le recuit
subs\'equent d\'eveloppe la texture cube
$\{100\}\langle 001\rangle$ comme composante de recristallisation
dominante.

\subsection{Le m\'ecanisme de Ridha--Hutchinson}

Ridha et Hutchinson \cite{Ridha1982} ont propos\'e un m\'ecanisme pour
expliquer pourquoi les r\'egions d'orientation cube dans la
microstructure d\'eform\'ee servent de sites de nucl\'eation
pr\'ef\'erentiels.  Lors du laminage en d\'eformation plane, quatre
syst\`emes de glissement sont actifs mais avec seulement deux vecteurs
de Burgers distincts.  Dans les bandes d'orientation cube, cette
activit\'e de glissement \'equilibr\'ee produit une sous-structure
relativement stable \`a faible densit\'e de dislocations.  Ces bandes
cube restaurent facilement pour former des sous-grains bord\'es de
joints \`a grand angle---les conditions pr\'ealables \`a des germes de
recristallisation viables.  Les germes r\'esultants croissent rapidement
dans la matrice de fibre~$\beta$ fortement d\'eform\'ee environnante,
expliquant la dominance de la composante cube.

Ce mod\`ele a \'et\'e largement cit\'e et reste l'explication standard
de la nucl\'eation de la texture cube dans les m\'etaux CFC lamin\'es et
recuits.

%% ===================================================================
\section{Aciers \'electriques \`a grains orient\'es}
\label{sec:goss}

Hutchinson a contribu\'e \`a la compr\'ehension de la recristallisation
secondaire et de la texture de Goss $\{110\}\langle 001\rangle$ dans
les aciers \'electriques Fe--Si.  Dans l'acier \'electrique \`a grains
orient\'es, une tr\`es faible fraction ($\sim\!10^{-6}$) des grains
primaires recristallis\'es d'orientation Goss cro\^it de mani\`ere
anormale lors d'un recuit \`a haute temp\'erature, consommant tous les
autres grains.  Cette \og recristallisation secondaire \fg{} produit une
texture de Goss aigu\"e qui conf\`ere d'excellentes propri\'et\'es
magn\'etiques (faibles pertes fer, haute perm\'eabilit\'e selon la
direction de laminage).

Hutchinson et Homma \cite{Homma2003} ont \'etudi\'e la d\'ependance
en orientation de la recristallisation secondaire dans le fer--silicium,
fournissant des preuves exp\'erimentales sur les d\'esorientations de
joints de grains qui favorisent la croissance anormale des grains et
s\'electionnent ainsi l'orientation de Goss.  Il a \'egalement examin\'e
la contribution des bandes de cisaillement \`a l'origine des germes
d'orientation Goss dans la texture de recristallisation primaire, reliant
ses travaux plus larges sur la sous-structure de d\'eformation \`a ce
probl\`eme d'importance industrielle critique.

%% ===================================================================
\section{Aciers microalli\'es et vanadium}
\label{sec:microalloy}

Au SIMR, Hutchinson a co-\'ecrit une monographie majeure sur le r\^ole
du vanadium dans les aciers microalli\'es \cite{Lagneborg2014}, avec
R.~Lagneborg, T.~Siwecki et S.~Zajac.

\subsection{M\'etallurgie physique du vanadium dans l'acier}

La monographie compare syst\'ematiquement le vanadium aux autres
\'el\'ements de microalliage (Nb, Ti) en termes de stabilit\'e
thermodynamique, solubilit\'e dans l'aust\'enite et comportement de
pr\'ecipitation.  Une distinction cl\'e est que V(C,N) pr\'esente une
solubilit\'e relativement \'elev\'ee dans l'aust\'enite compar\'ee \`a
NbC ou TiN.  Aux temp\'eratures typiques de fin de laminage
($>900\,{}^\circ$C), la quasi-totalit\'e du vanadium reste en solution
solide, alors que le Nb pr\'ecipite abondamment dans l'aust\'enite \`a
ces temp\'eratures.  En cons\'equence, la contribution principale du
vanadium au renforcement se fait par durcissement par pr\'ecipitation
dans la ferrite, plut\^ot que par affinement des grains
aust\'enitiques (qui est le domaine du Nb et du Ti).

\subsection{M\'ecanismes de pr\'ecipitation dans la ferrite}

La monographie identifie trois modes distincts de pr\'ecipitation de
V(C,N) dans la ferrite~:

\begin{enumerate}
  \item \textbf{La pr\'ecipitation interphase} se produit \`a
    l'interface $\gamma/\alpha$ en migration lors de la transformation
    aust\'enite--ferrite.  Des rang\'ees de pr\'ecipit\'es se forment
    aux positions successives de l'interface, produisant les alignements
    planaires caract\'eristiques observ\'es en MET.  Lagneborg et Zajac
    \cite{LagneborgZajac2001} ont propos\'e un mod\`ele de migration par
    marches dans lequel V(C,N) nucl\'ee sur les terrasses stationnaires
    de l'interface~; lorsque les marches balaient la surface, la terrasse
    charg\'ee de pr\'ecipit\'es est laiss\'ee en arri\`ere et une
    nouvelle terrasse est expos\'ee pour la g\'en\'eration suivante de
    nucl\'eation.  L'espacement entre feuillets est r\'egi par la
    hauteur des marches et la temp\'erature de transformation.

  \item \textbf{La pr\'ecipitation g\'en\'erale (al\'eatoire)} se
    produit apr\`es l'ach\`evement de la transformation, lors du
    refroidissement lent ou du maintien isotherme dans la ferrite.
    Contrairement \`a la pr\'ecipitation interphase, les particules
    nucl\'eent sur les dislocations et les sous-joints, produisant une
    distribution spatiale al\'eatoire.

  \item \textbf{La pr\'ecipitation fibreuse} peut se produire lorsque
    l'interface se d\'eplace tr\`es lentement, produisant des fibres de
    pr\'ecipit\'es allong\'es plut\^ot que des particules discr\`etes.
\end{enumerate}

\`A partir de la relation d'Orowan, la monographie d\'erive des
estimations quantitatives de renforcement~: environ 6\,MPa par
0.001\,\%\,N en masse et 5,5\,MPa par 0,01\,\%\,C en masse pour les
V(C,N) pr\'ecipit\'es en interphase.  Un r\'esultat important est que
le carbone et l'azote contribuent de mani\`ere synergique au
durcissement par pr\'ecipitation.  \`A des teneurs plus \'elev\'ees en
azote, les particules de V(C,N) sont plus fines et plus densement
distribu\'ees parce que VN a un produit de solubilit\'e plus faible que
VC, conduisant \`a un taux de nucl\'eation plus \'elev\'e.  La
dispersion plus fine incorpore alors \'egalement davantage de carbone,
de sorte que les aciers combin\'es V--N--C atteignent un durcissement
par pr\'ecipitation sup\'erieur \`a la somme des contributions V--N et
V--C prises s\'epar\'ement.

\subsection{VN comme germinateur de ferrite}

La monographie d\'ecrit un m\'ecanisme d'affinement des grains de
ferrite par nucl\'eation intragranulaire sur des pr\'ecipit\'es de VN
dans l'aust\'enite.  Les particules de VN qui se forment lors du
refroidissement \`a travers le domaine aust\'enitique (ou lors d'une
\'etape de transfert/r\'echauffage) adoptent la relation d'orientation
de Baker--Nutting avec la ferrite,
$(100)_\text{VN}\parallel(100)_\alpha$,
$[011]_\text{VN}\parallel[001]_\alpha$.  Cette interface \`a faible
\'energie fait du VN un site efficace de nucl\'eation h\'et\'erog\`ene
pour la ferrite, favorisant la formation de ferrite intragranulaire et
produisant un rapport de transformation $D_\gamma/D_\alpha$ aussi
\'elev\'e que 5--10 (contre $\sim$2 pour la nucl\'eation conventionnelle
aux joints de grains).  Les particules de TiN peuvent servir de
pr\'ecurseurs \`a la pr\'ecipitation de VN par co-pr\'ecipitation,
renfor\c{c}ant ce m\'ecanisme dans les aciers Ti--V--N.

\subsection{Laminage contr\^ol\'e par recristallisation (RCR)}

La haute solubilit\'e de V(C,N) dans l'aust\'enite signifie que le
vanadium ne retarde pas la recristallisation entre les passes de
laminage dans la m\^eme mesure que le niobium.  En cons\'equence, le
microalliage au vanadium ne se pr\^ete pas au laminage contr\^ol\'e
conventionnel (CR) en dessous de la temp\'erature d'arr\^et de la
recristallisation.  Au lieu de cela, Hutchinson et ses collaborateurs ont
d\'evelopp\'e le concept de \emph{laminage contr\^ol\'e par
recristallisation} (RCR) \cite{Zajac1991,Siwecki1995}~:

\begin{itemize}
  \item Le laminage est effectu\'e \`a des temp\'eratures de finition
    relativement \'elev\'ees ($\geq 950\,{}^\circ$C), o\`u l'aust\'enite
    recristallise compl\`etement entre les passes, maintenant une
    structure granulaire fine \'equiaxe par des cycles de
    recristallisation r\'ep\'et\'es.
  \item De faibles additions de Ti ($\sim$0,01\,\%) fournissent des
    particules de TiN qui \'epinglent les joints de grains
    aust\'enitiques par tra\^in\'ee de Zener, emp\^echant le
    grossissement des grains apr\`es recristallisation.
  \item Un refroidissement acc\'el\'er\'e (ACC) apr\`es laminage, suivi
    d'un bobinage \`a 550--$600\,{}^\circ$C pour les bandes, entra\^ine
    la transformation $\gamma\to\alpha$ vers des microstructures fines de
    ferrite--bainite et pr\'eserve le vanadium en solution pour une
    pr\'ecipitation fine de V(C,N) lors du refroidissement lent dans
    la bobine.
\end{itemize}

Le mod\`ele informatique MICDEL, d\'evelopp\'e \`a l'Institut su\'edois
de recherche sur les m\'etaux, a \'et\'e utilis\'e pour optimiser les
programmes de passes RCR en simulant l'\'evolution coupl\'ee de la
taille des grains d'aust\'enite, de la fraction recristallis\'ee et de
la distribution des pr\'ecipit\'es pendant le laminage \`a chaud
\cite{Siwecki1992}.  Un r\'esultat pratique critique \'etait que les
petites passes de \og calibrage \fg{} en fin de s\'equence de laminage
doivent \^etre \'evit\'ees~: de faibles d\'eformations \`a haute
temp\'erature peuvent provoquer un grossissement anormal des grains
plut\^ot qu'un affinement, parce que la nucl\'eation de la
recristallisation est insuffisamment abondante (le ph\'enom\`ene de
\og recuit critique \fg{}).

\subsection{Effet de l'azote}

Un th\`eme central de la monographie est le r\^ole puissant et souvent
sous-estim\'e de l'azote dans les aciers microalli\'es au V.
L'augmentation de l'azote de $\sim$0,003\,\% \`a $\sim$0,013\,\% dans
un acier \`a 0,09\,\%V augmente la limite d'\'elasticit\'e de
$\sim$120\,MPa pour un mat\'eriau trait\'e par RCR \cite{Zajac1991}.
Cela provient de deux m\'ecanismes qui se renforcent mutuellement~:
(i)~une nucl\'eation plus dense de VN produit des dispersions de
pr\'ecipit\'es plus fines avec un espacement inter-particules plus
faible, am\'eliorant le durcissement d'Orowan~; (ii)~les particules de
VN dans l'aust\'enite servent de germinateurs intragranulaires de
ferrite (Section~\ref{sec:microalloy}), affinant la taille des grains de
ferrite.  La monographie d\'emontre que l'azote est un \'el\'ement de
microalliage essentiel en conjonction avec le vanadium---et non
simplement une impuret\'e d\'el\'et\`ere comme on le consid\'erait
traditionnellement.

%% ===================================================================
\section{T\^ole d'aluminium~: texture, cornes et formabilit\'e}
\label{sec:aluminium}

Hutchinson a appliqu\'e son expertise en texture \`a la mise en forme des
t\^oles d'aluminium.  Les travaux avec Oscarsson et Ekstr\"om
\cite{Oscarsson1991} ont \'etudi\'e l'influence de la microstructure
initiale sur la texture et les cornes dans la t\^ole d'aluminium apr\`es
laminage \`a froid et recuit.

\subsection{Contr\^ole des cornes}

Les \og cornes \fg{}---la formation d'un bord ondul\'e (oreilles) lors
de l'emboutissage profond d'un flan circulaire---sont une manifestation
directe de l'anisotropie planaire caus\'ee par la texture
cristallographique.  Une texture de laminage produit des cornes \`a
$0^\circ/90^\circ$~; une texture de recristallisation cube produit des
cornes \`a $45^\circ$.  En \'equilibrant les intensit\'es relatives de
ces deux composantes de texture par le contr\^ole de la temp\'erature de
fin de laminage \`a chaud, du degr\'e de recristallisation avant le
laminage \`a froid, du taux de r\'eduction au laminage \`a froid et des
conditions de recuit, les cornes peuvent \^etre minimis\'ees.

\subsection{Mod\'elisation de la texture \`a travers le proc\'ed\'e}

La compr\'ehension de Hutchinson des m\'ecanismes d'\'evolution de la
texture---formation de la texture de d\'eformation, restauration,
nucl\'eation et croissance des grains---a fourni la base physique des
mod\`eles de texture \`a travers le proc\'ed\'e appliqu\'es \`a la
production de t\^oles d'aluminium, reliant les param\`etres de
proc\'ed\'e \`a la texture finale et, de l\`a, \`a l'anisotropie
m\'ecanique, au comportement des cornes et \`a la formabilit\'e.

%% ===================================================================
\section{Ultrasons laser pour le contr\^ole en ligne des proc\'ed\'es}
\label{sec:lus}

Dans des travaux plus r\'ecents, Hutchinson a \'etudi\'e les techniques
d'ultrasons laser (LUS) pour la mesure sans contact et en temps r\'eel
de la microstructure pendant le laminage \`a chaud de l'acier
\cite{Hutchinson2011LUS}.  Des essais au laminoir \`a bandes \`a chaud
de SSAB \`a Borl\"ange, en Su\`ede, ont d\'emontr\'e que des ondes
ultrasonores g\'en\'er\'ees et d\'etect\'ees par laser peuvent mesurer
l'\'etat de transformation aust\'enite--ferrite, la fraction
recristallis\'ee et la taille des grains sur une bande en mouvement \`a
la table de sortie.  Ces travaux repr\'esentent une fronti\`ere dans le
contr\^ole en boucle ferm\'ee de la microstructure pour la production
industrielle d'acier, reliant la compr\'ehension approfondie de
Hutchinson des transformations de phase et de la cin\'etique de
recristallisation au suivi pratique des proc\'ed\'es.

%% ===================================================================
\section{Microstructures de d\'eformation et textures de transformation}
\label{sec:deformation}

Un article de 1999 dans \emph{Philosophical Transactions of the Royal
Society} \cite{Hutchinson1999Phil} a fourni un traitement complet des
microstructures et des textures de d\'eformation dans les aciers.  Ce
travail a unifi\'e les observations de Hutchinson sur la mani\`ere dont
diff\'erents modes de d\'eformation---d\'eformation plane (laminage),
traction uniaxiale, cisaillement---cr\'eent des sous-structures de
dislocations caract\'eristiques dans le fer CC, et comment ces
sous-structures contr\^olent \`a leur tour les textures qui se
d\'eveloppent lors de la recristallisation ou de la transformation de
phase ult\'erieure.  Avec ses co-auteurs Ryde et Bate, il a \'egalement
\'etudi\'e les textures de transformation issues de la transformation
aust\'enite--ferrite, examinant comment les relations d'orientation de
Kurdjumov--Sachs et de Nishiyama--Wassermann entre l'aust\'enite
parente et la ferrite produite influencent la texture de la
microstructure ferritique finale.

%% ===================================================================
\section{Anisotropie dans les m\'etaux~: une synth\`ese}
\label{sec:anisotropy}

Une \'evaluation critique de 2015 dans \emph{Materials Science and
Technology}, \og Anisotropy in metals \fg{} \cite{Hutchinson2015}, sert
de synth\`ese de l'ensemble des th\`emes de la carri\`ere de Hutchinson.
Elle traite de l'anisotropie \'elastique et plastique dans les
syst\`emes CC et CFC, montrant comment la texture cristallographique
contr\^ole les propri\'et\'es m\'ecaniques directionnelles.
L'\'evaluation couvre les constantes \'elastiques du monocristal, les
sch\'emas de moyenne polycristalline (Voigt, Reuss, Hill, Kr\"oner), la
relation entre la FDO (fonction de distribution des orientations) et les
valeurs de $r$, et les implications pratiques pour la mise en forme des
t\^oles.  Elle d\'emontre le fil conducteur des travaux de Hutchinson~:
la texture est le pont entre la mise en \oe uvre, la microstructure et
les performances m\'ecaniques.

Dans un article de 2025 \cite{Hutchinson2025pssb}, Hutchinson a
poursuivi ce th\`eme en mesurant exp\'erimentalement les vitesses des
ondes \'elastiques longitudinales dans le fer polycristallin de la
temp\'erature ambiante \`a $880\,{}^\circ\mathrm{C}$, comparant les
r\'esultats avec les mod\`eles de moyenne de Hill et Kr\"oner.

%% ===================================================================
\section{Autres travaux notables}
\label{sec:other}

\subsection{Manuel~: An Introduction to Textures in Metals}

Hutchinson a co-\'ecrit avec M.~Hatherly la monographie \emph{An
Introduction to Textures in Metals} (Institute of Metals, Londres,
1979~; 2\ieme{}~\'ed.\ 1990) \cite{Hatherly1979}.  Ce manuel a servi de
r\'ef\'erence standard dans le domaine de la texture cristallographique
des m\'etaux.

\subsection{La m\'et\'eorite de Gibeon}

Dans une application inventive des techniques d'EBSD (diffraction des
\'electrons r\'etrodiffus\'es), Hutchinson et Hagstr\"om
\cite{Hutchinson2006Met} ont \'etudi\'e la transformation ferritique de
Widmanst\"atten dans la m\'et\'eorite de fer de Gibeon, trouvant des
relations d'orientation entre les plaquettes de ferrite et l'aust\'enite
parente conformes aux variantes de Nishiyama--Wassermann et de
Kurdjumov--Sachs---une exp\'erience naturelle de transformation de phase
survenue sur des \'echelles de temps g\'eologiques.

\subsection{Recuit partiel des aciers lamin\'es \`a froid}

Avec R.K.~Ray et C.~Ghosh \cite{Ray2011}, Hutchinson a pass\'e en revue
le recuit partiel (\emph{back-annealing})---un traitement thermique des
aciers lamin\'es \`a froid con\c{c}u pour obtenir une restauration
partielle et/ou une recristallisation partielle plut\^ot qu'une
recristallisation compl\`ete.  Cette voie de proc\'ed\'e permet
d'ajuster le compromis r\'esistance--ductilit\'e en contr\^olant la
fraction de microstructure restaur\'ee par rapport \`a celle
recristallis\'ee.

%% ===================================================================
\section{La monographie sur le vanadium~: un r\'esum\'e complet}
\label{sec:vanadium-monograph}

L'\oe uvre unique la plus substantielle dans laquelle la pens\'ee
m\'etallurgique de Hutchinson est expos\'ee est la monographie de
101~pages \emph{The Role of Vanadium in Microalloyed Steels}
\cite{Lagneborg2014}, r\'edig\'ee avec R.~Lagneborg, T.~Siwecki et
S.~Zajac \`a l'Institut su\'edois de recherche sur les m\'etaux
(aujourd'hui Swerim).  Publi\'ee initialement en 1999 dans le
\emph{Scandinavian Journal of Metallurgy}, une \'edition \'elargie a
paru sous la r\'ef\'erence Swerea KIMAB report KIMAB-2014-115, avec le
soutien du Comit\'e technique international du vanadium (Vanitec).  La
monographie est d\'edi\'ee \`a la m\'emoire de Michael Korchynsky,
pionnier du microalliage au vanadium.  Ce qui suit est un r\'esum\'e
autonome de cette \oe uvre, r\'edig\'e pour un lecteur scientifiquement
form\'e qui n'est pas n\'ecessairement sp\'ecialiste en m\'etallurgie.

\subsection{Contexte~: qu'est-ce qu'un acier microalli\'e~?}

Les aciers de construction---utilis\'es dans les b\^atiments, les ponts,
les navires, les pipelines et les v\'ehicules---doivent combiner haute
r\'esistance avec bonne ductilit\'e et t\'enacit\'e (r\'esistance \`a la
rupture fragile et soudaine).  La mani\`ere traditionnelle de renforcer
l'acier est d'ajouter davantage de carbone, mais le carbone rend l'acier
fragile et difficile \`a souder.  Les aciers \emph{microalli\'es} (ou
HSLA, \emph{High-Strength Low-Alloy}) r\'esolvent ce dilemme en ajoutant
de tr\`es petites quantit\'es---typiquement 0,01--0,15\,\%---d'\'el\'ements
tels que le vanadium~(V), le niobium~(Nb) et le titane~(Ti).  Ces
\og micro-additions \fg{} r\'ealisent leur effet non par renforcement en
solution solide (qui n\'ecessite de grandes quantit\'es) mais par deux
m\'ecanismes plus efficaces~:
\begin{enumerate}
  \item \textbf{Affinement des grains~:} rendre les grains cristallins de
    l'acier plus fins, ce qui augmente simultan\'ement la r\'esistance et
    la t\'enacit\'e---une combinaison rare en science des mat\'eriaux,
    puisque ces propri\'et\'es sont habituellement antag\-onistes.
  \item \textbf{Durcissement par pr\'ecipitation~:} former un nombre
    consid\'erable de particules de carbures et de nitrures \`a
    l'\'echelle nanom\'etrique qui entravent le mouvement des dislocations
    (les d\'efauts lin\'eaires qui portent la d\'eformation plastique),
    \'elevant la limite d'\'elasticit\'e.
\end{enumerate}

Les trois \'el\'ements de microalliage (V, Nb, Ti) diff\`erent
fondamentalement par le \emph{moment} et le \emph{lieu} o\`u leurs
particules se forment, et la monographie est structur\'ee autour de
cette distinction.

\subsection{La base thermodynamique~: la solubilit\'e compte}

L'acier est mis en forme par chauffage \`a haute temp\'erature o\`u il
existe dans la structure cristalline cubique \`a faces centr\'ees (CFC)
appel\'ee \emph{aust\'enite} ($\gamma$), puis lamin\'e \`a chaud, et
enfin refroidi pour qu'il se transforme en la phase cubique centr\'ee
(CC) appel\'ee \emph{ferrite} ($\alpha$) qui constitue le produit final.
La question cruciale pour chaque \'el\'ement de microalliage est~: quelle
quantit\'e se dissout dans l'aust\'enite \`a la temp\'erature de
laminage, et quelle quantit\'e a d\'ej\`a pr\'ecipit\'e~?

\begin{itemize}
  \item \textbf{TiN} est extr\^emement stable et presque insoluble dans
    l'aust\'enite.  Il pr\'ecipite lors de la coul\'ee et persiste \`a
    travers toutes les \'etapes ult\'erieures de mise en \oe uvre.  Son
    r\^ole est d'\'epingler les joints de grains aust\'enitiques
    (\'epinglage de Zener), emp\^echant le grossissement des grains lors
    du r\'echauffage et du soudage.
  \item \textbf{NbC} a une stabilit\'e interm\'ediaire.  Il pr\'ecipite
    dans l'aust\'enite pendant le laminage \`a chaud, en particulier en
    dessous de $\sim$950\,${}^\circ$C.  Le Nb en solution solide et sous
    forme de pr\'ecipit\'es retardent tous deux fortement la
    recristallisation de l'aust\'enite, produisant des grains
    d'aust\'enite allong\'es (\og pancakes \fg{}) qui se transforment en
    ferrite tr\`es fine.  Cependant, la majeure partie du Nb pr\'ecipite
    dans l'aust\'enite, de sorte que relativement peu reste disponible
    pour le durcissement par pr\'ecipitation dans la ferrite.
  \item \textbf{V(C,N)} a une haute solubilit\'e dans l'aust\'enite.
    Aux temp\'eratures typiques de fin de laminage, la quasi-totalit\'e
    du vanadium reste dissous.  Il ne pr\'ecipite que pendant ou apr\`es
    la transformation $\gamma\to\alpha$, formant des particules de V(C,N)
    \`a l'\'echelle nanom\'etrique dans la ferrite qui assurent un fort
    durcissement par pr\'ecipitation.
\end{itemize}

Ces diff\'erences ne sont pas des d\'efauts---elles sont
compl\'ementaires.  La monographie soutient que V, Nb et Ti doivent
\^etre consid\'er\'es comme remplissant des r\^oles diff\'erents dans
une conception d'alliage coordonn\'ee.

\subsection{Trois modes de pr\'ecipitation dans la ferrite}

Lorsque l'aust\'enite se transforme en ferrite lors du refroidissement,
le vanadium qui \'etait dissous dans l'aust\'enite devient sursatur\'e
dans la ferrite et pr\'ecipite.  La monographie identifie trois modes~:

\begin{enumerate}
  \item \textbf{Pr\'ecipitation interphase.}  Lorsque le front de
    transformation $\gamma/\alpha$ balaye l'acier, des rang\'ees de
    particules de V(C,N) nucl\'eent sur l'interface elle-m\^eme.
    L'interface avance par un m\'ecanisme de marches~: les terrasses
    planes accumulent les pr\'ecipit\'es, tandis que les marches
    lat\'erales balayent la surface, laissant la terrasse charg\'ee de
    pr\'ecipit\'es en arri\`ere et exposant une surface fra\^iche pour
    la rang\'ee suivante.  En microscopie \'electronique en transmission,
    cela produit des alignements saisissants de feuillets de
    pr\'ecipit\'es parall\`eles, espac\'es de 10--30\,nm.

  \item \textbf{Pr\'ecipitation g\'en\'erale (al\'eatoire).}  Apr\`es le
    passage du front de transformation, la sursaturation r\'esiduelle
    entra\^ine une pr\'ecipitation suppl\'ementaire sur les dislocations
    et les sous-joints dans la ferrite.  Ces particules sont distribu\'ees
    de mani\`ere al\'eatoire plut\^ot qu'arrang\'ees en feuillets.

  \item \textbf{Pr\'ecipitation fibreuse.}  \`A des vitesses d'interface
    tr\`es faibles, les pr\'ecipit\'es peuvent cro\^itre sous forme de
    fibres continues plut\^ot que de particules discr\`etes.  Ce cas est
    moins courant dans les conditions industrielles.
\end{enumerate}

Les pr\'ecipitations interphase et g\'en\'erale contribuent toutes deux
\`a l'incr\'ement de durcissement d'Orowan (l'augmentation de la limite
d'\'elasticit\'e caus\'ee par les particules que les dislocations doivent
contourner).  La monographie donne des estimations quantitatives~:
environ 6\,MPa par 0,001\,\%\,N en masse et 5,5\,MPa par 0,01\,\%\,C
en masse.

\subsection{L'azote~: d'impuret\'e \`a \'el\'ement d'alliage}

L'un des arguments les plus frappants de la monographie est la
r\'ehabilitation de l'azote en tant qu'\'el\'ement d'alliage
\emph{b\'en\'efique}.  Conventionnellement, l'azote dans l'acier est
consid\'er\'e comme nocif~: l'azote \og libre \fg{} provoque le
vieillissement sous contrainte et la fragilisation.  Cependant, dans les
aciers contenant du vanadium, l'azote est li\'e sous forme de VN, qui
est thermodynamiquement plus stable que VC.  Le taux de nucl\'eation
plus \'elev\'e de VN produit des dispersions de pr\'ecipit\'es plus
fines et plus denses, am\'eliorant consid\'erablement le renforcement.
La monographie montre que le doublement de la teneur en azote dans un
acier pour armatures \`a 0,11\,\%V (de $\sim$85\,ppm \`a
$\sim$180\,ppm) augmente la limite d'\'elasticit\'e de pr\`es de
120\,MPa---un effet remarquable pour un \'el\'ement \`a l'\'etat de
traces.

De plus, le carbone et l'azote agissent de mani\`ere synergique.  Les
fines particules de VN qui nucl\'eent en premier servent de germes pour
la croissance ult\'erieure de V(C,N) incorporant du carbone.  Le
renforcement combin\'e d\'epasse ce que chaque \'el\'ement fournirait
seul.

\subsection{Nucl\'eation de la ferrite sur VN~: transformer la
  transformation de l'aust\'enite en outil d'affinement}

Une contribution particuli\`erement \'el\'egante est le concept
d'utilisation des pr\'ecipit\'es de VN comme germinateurs
intragranulaires de ferrite.  Normalement, les grains de ferrite
nucl\'eent aux joints de grains aust\'enitiques lors du
refroidissement~; la taille des grains de ferrite r\'esultante est
limit\'ee par la taille des grains d'aust\'enite ant\'erieurs.  Mais les
particules de VN qui pr\'ecipitent \`a l'int\'erieur des grains
d'aust\'enite (lors d'un cycle de refroidissement/r\'echauffage ou d'un
maintien prolong\'e) ont une relation d'orientation cristallographique
avec la ferrite (la relation de Baker--Nutting) qui en fait des sites de
nucl\'eation h\'et\'erog\`ene puissants.  Les grains de ferrite
nucl\'eent \`a l'int\'erieur des grains d'aust\'enite sur ces
particules, produisant un rapport de raffinement par transformation
$D_\gamma/D_\alpha$ de 5--10 (contre $\sim$2 pour la nucl\'eation
conventionnelle aux joints de grains seuls).  Des essais industriels
chez Riva Acciaio Sellero (Italie) sur des poutrelles lourdes ont
confirm\'e des r\'eductions de taille de grains de 15\,$\mu$m \`a
10\,$\mu$m et des augmentations de la limite d'\'elasticit\'e de
$\sim$130\,MPa.

\subsection{Voies de mise en \oe uvre~: RCR, CR et le r\^ole du
  refroidissement}

La monographie consacre une attention consid\'erable \`a la mani\`ere
dont les aciers au vanadium doivent \^etre lamin\'es \`a chaud.  Il
existe deux voies principales~:

\paragraph{Laminage contr\^ol\'e (CR).} L'approche conventionnelle~:
laminage de finition \`a basse temp\'erature (en dessous de
$\sim$900\,${}^\circ$C) pour d\'eformer l'aust\'enite sans
recristallisation, cr\'eant des grains allong\'es en \og pancakes \fg{}
dont la grande surface de joints produit de la ferrite fine.  Cela
n\'ecessite des efforts de laminage \'elev\'es et convient aux aciers au
Nb, qui r\'esistent fortement \`a la recristallisation.

\paragraph{Laminage contr\^ol\'e par recristallisation (RCR).}
L'approche favoris\'ee pour les aciers au vanadium~: laminage de
finition \`a des temp\'eratures plus \'elev\'ees
($>$\,950\,${}^\circ$C), o\`u l'aust\'enite recristallise compl\`etement
entre les passes.  L'affinement des grains provient de la
recristallisation r\'ep\'et\'ee plut\^ot que du \og pancaking \fg{}.
Les particules de TiN emp\^echent le grossissement des grains entre les
passes.  Le RCR n\'ecessite des efforts de laminage plus faibles
(environ 25\,\% de moins que le CR pour le m\^eme produit), offre une
productivit\'e du laminoir sup\'erieure, et constitue la voie naturelle
pour le vanadium car la haute solubilit\'e de V(C,N) signifie qu'il n'y
a pas de pr\'ecipit\'es pour \'epingler l'aust\'enite contre la
recristallisation.

Apr\`es le laminage, un \emph{refroidissement acc\'el\'er\'e} sur la
table de sortie transforme l'aust\'enite fine en un m\'elange encore plus
fin de ferrite--bainite.  Dans les produits en bande, la bande est
ensuite bobin\'ee \`a 550--600\,${}^\circ$C, o\`u le refroidissement lent
dans la bobine fournit des conditions id\'eales pour la pr\'ecipitation
de V(C,N).  La monographie montre que cette voie RCR+ACC atteint des
niveaux de r\'esistance et de t\'enacit\'e comparables ou sup\'erieurs au
CR conventionnel, tout en \'etant industriellement plus efficace.

\subsection{Applications industrielles}

Les chapitres~7 et~8 de la monographie passent en revue des applications
produit sp\'ecifiques~:

\paragraph{T\^oles fortes.}  Des essais \`a l'\'echelle industrielle sur
des t\^oles \`a 0,01\,\%Ti--0,08\,\%V \`a Pohang (Cor\'ee),
Rautaruukki (Finlande) et Baoshan (Chine) ont atteint des limites
d'\'elasticit\'e sup\'erieures \`a 490\,MPa avec une excellente
t\'enacit\'e aux chocs (ITT$_{40\text{J}} < -80\,{}^\circ$C) par un
traitement RCR+ACC.

\paragraph{Bandes lamin\'ees \`a chaud.}  Le laminage de bandes est
presque id\'ealement adapt\'e au microalliage V--N~: le refroidissement
rapide \`a l'eau sur la table de sortie pr\'eserve le vanadium en
solution, tandis que le refroidissement lent dans la bobine fournit
amplement de temps pour une pr\'ecipitation fine de V(C,N).  Des limites
d'\'elasticit\'e de 550\,MPa ont \'et\'e atteintes dans des bandes de
8--10\,mm avec des temp\'eratures de transition aux chocs inf\'erieures
\`a $-40\,{}^\circ$C.

\paragraph{Bandes bainitiques.}  Un d\'eveloppement notable chez Swerea
KIMAB et SSAB fut un acier bainitique microalli\'e au V visant une
limite d'\'elasticit\'e de 700\,MPa pour des applications de ch\^assis
de camion \cite{Hutchinson2014Bainite}.  Dans ce syst\`eme, le vanadium
ne modifie pas la transformation bainitique elle-m\^eme mais inhibe la
restauration de la structure de dislocations pendant le bobinage,
pr\'eservant la r\'esistance m\^eme \`a des temp\'eratures de bobinage
inf\'erieures \`a 500\,${}^\circ$C---un avantage pratique \'etant donn\'e
la difficult\'e du contr\^ole pr\'ecis de la temp\'erature de bobinage
sur les tables de sortie industrielles.

\paragraph{Armatures (ronds \`a b\'eton).}  Les armatures constituent la
plus grande application du microalliage au V en tonnage.  Avec des
teneurs en carbone plus \'elev\'ees ($\sim$0,16--0,22\,\%) et des
temp\'eratures de finition \'elev\'ees, tout le vanadium reste en
solution pour le durcissement par pr\'ecipitation.  La monographie
montre que 5~parties d'N en masse \'equivalent \`a 1~partie de V en
termes de renforcement.  Le microalliage V--N est devenu la voie
privil\'egi\'ee pour les armatures \`a haute r\'esistance hors d'Europe,
avec une production particuli\`erement importante en Chine
($\sim$130\,Mton/an d'armatures en 2010).

\paragraph{Pi\`eces forg\'ees.}  Les aciers de forge microalli\'es au V
\`a structure ferrite--perlite (par exemple 30MNV56) offrent une
alternative comp\'etitive aux nuances tremp\'ees et revenues (Q-T) pour
les composants automobiles tels que les vilebrequins et les bielles.
Ils \'eliminent les \'etapes de traitement thermique (trempe, revenu,
dressage), r\'eduisent les co\^uts d'usinage de 28\,\% et pr\'esentent
une r\'esistance \`a la fatigue comparable, quoique avec une t\'enacit\'e
quelque peu inf\'erieure.

\paragraph{Tubes sans soudure.}  Une application innovante exploite la
nucl\'eation intragranulaire de ferrite sur VN pour \'eliminer l'\'etape
de normalisation en ligne traditionnellement requise pour affiner la
microstructure grossi\`ere produite par le per\c{c}age et le laminage \`a
haute temp\'erature.

\subsection{Soudabilit\'e}

Le chapitre~8 aborde une pr\'eoccupation cruciale~: la teneur plus
\'elev\'ee en azote b\'en\'efique pour la r\'esistance compromet-elle la
soudabilit\'e~?  La zone affect\'ee thermiquement (ZAT) adjacente \`a
une soudure est r\'echauff\'ee \`a proximit\'e du point de fusion,
dissolvant les pr\'ecipit\'es et grossissant la structure des grains
d'aust\'enite.  La monographie montre que~:

\begin{itemize}
  \item De faibles additions de Ti ($\sim$0,01\,\%) sont essentielles~:
    les particules de TiN r\'esistent aux temp\'eratures de soudage et
    limitent le grossissement des grains d'aust\'enite dans la ZAT,
    maintenant la t\'enacit\'e \`a la ligne de fusion.
  \item Pour le soudage \`a \'energie faible \`a moyenne
    ($\Delta t_{8/5} < 30$\,s), les aciers \`a plus haute teneur en azote
    pr\'esentent en fait une \emph{meilleure} t\'enacit\'e de la ZAT
    \`a gros grains que les variantes \`a faible azote, parce que le V et
    le N en solution favorisent la ferrite aciculaire plut\^ot que la
    ferrite grossi\`ere aux joints de grains, produisant une taille de
    grain effective plus fine.
  \item Pour les fortes \'energies de soudage ($>$10\,kJ/mm), la
    t\'enacit\'e peut se d\'et\'eriorer dans les aciers \`a haute teneur
    en N en raison de la ferrite polygonale grossi\`ere aux joints de
    grains~; ceci est att\'enu\'e par des micro-additions de Ti/REM/Ca.
  \item Le soudage multipasse \emph{am\'eliore} en fait la t\'enacit\'e
    de la ZAT~: le VN dissous par la premi\`ere passe re-pr\'ecipite
    pendant le cycle thermique des passes ult\'erieures, affinant la
    microstructure.
\end{itemize}

La conclusion g\'en\'erale est que les aciers microalli\'es V--N sont
pleinement compatibles avec une bonne soudabilit\'e lorsque des limites
raisonnables sur l'\'energie de soudage sont respect\'ees et que de
faibles additions de Ti sont incluses.

\subsection{La monographie en contexte}

La monographie de Lagneborg--Hutchinson--Siwecki--Zajac occupe une
position unique dans la litt\'erature sur les aciers microalli\'es.  Ses
plus proches \'equivalents sont l'ouvrage de Gladman \emph{The Physical
Metallurgy of Microalloyed Steels} (1997), qui offre un panorama plus
large mais moins d\'etaill\'e de tous les \'el\'ements de microalliage,
et les actes des conf\'erences Microalloying (1975, 1995, 2007) qui
rassemblent des contributions plus courtes de nombreux groupes de
recherche.  La monographie sur le vanadium se distingue par~:
\begin{itemize}
  \item Son approche syst\'ematique centr\'ee sur un \'el\'ement~: chaque
    chapitre est organis\'e autour de la question \og que fait
    sp\'ecifiquement le vanadium, et en quoi cela diff\`ere-t-il du Nb
    et du Ti~? \fg{}
  \item Son int\'egration de la thermodynamique fondamentale et de la
    th\'eorie de la pr\'ecipitation avec des donn\'ees industrielles
    grandeur nature provenant d'aci\'eries su\'edoises, cor\'eennes,
    finlandaises, italiennes, chinoises et br\'esiliennes.
  \item Son plaidoyer pour l'azote comme \'el\'ement d'alliage ajout\'e
    d\'elib\'er\'ement---une position qui \'etait controvers\'ee
    lorsqu'elle a \'et\'e avanc\'ee par le groupe Swerea dans les
    ann\'ees 1990 mais qui a depuis \'et\'e largement valid\'ee
    industriellement.
  \item Son d\'eveloppement de mod\`eles quantitatifs
    proc\'ed\'e--propri\'et\'es (MICDEL) qui comblent le foss\'e entre la
    compr\'ehension en laboratoire et la pratique industrielle.
\end{itemize}

La monographie repr\'esente l'aboutissement d'environ 30 ann\'ees de
travail collaboratif du groupe de microalliage du
SIMR/Swerea KIMAB, dont Hutchinson a \'et\'e un membre central tout au
long.

%% ===================================================================
\section*{R\'esum\'e}

Les contributions de Hutchinson peuvent \^etre regroup\'ees en plusieurs
th\`emes interconnect\'es~:

\begin{enumerate}
  \item \textbf{Relations texture--propri\'et\'es~:} \'etablir des liens
    quantitatifs entre la texture cristallographique (en particulier la
    fibre~$\gamma$ dans les CC et la composante cube dans les CFC) et
    l'anisotropie m\'ecanique macroscopique (valeurs de $\bar{r}$,
    cornes).

  \item \textbf{M\'ecanismes de recristallisation~:} identifier les
    sous-structures de d\'eformation (bandes de cisaillement, bandes de
    transition, bandes cube) qui servent de sites de nucl\'eation, et
    clarifier les r\^oles de la nucl\'eation orient\'ee et de la
    croissance s\'elective.

  \item \textbf{Contr\^ole proc\'ed\'e--texture~:} traduire la
    compr\'ehension m\'ecaniste en orientations pratiques pour la mise en
    \oe uvre industrielle---programmes de laminage \`a chaud, laminage
    \`a froid et recuit---afin d'atteindre les textures cibles dans les
    aciers et les alliages d'aluminium.

  \item \textbf{Aciers microalli\'es~:} \'elucider les r\^oles distincts
    de V, Nb et Ti dans la mise en forme thermom\'ecanique, avec des
    contributions particuli\`eres \`a la compr\'ehension de la
    pr\'ecipitation du vanadium et du laminage contr\^ol\'e par
    recristallisation.

  \item \textbf{Aciers \'electriques~:} contribuer \`a la compr\'ehension
    de la s\'election de la texture de Goss lors de la recristallisation
    secondaire dans les aciers Fe--Si \`a grains orient\'es.

  \item \textbf{Surveillance des proc\'ed\'es~:} \^etre pionnier dans
    l'application des ultrasons laser pour la mesure en ligne de la
    microstructure pendant le laminage \`a chaud des bandes.
\end{enumerate}

%% ===================================================================
\begin{thebibliography}{99}

\bibitem{Hutchinson1984IMR}
W.B.~Hutchinson,
``Development and control of annealing textures in low-carbon steels,''
\emph{International Materials Reviews}, vol.~29, no.~1, pp.~25--72, 1984.

\bibitem{Hutchinson2007MSF}
W.B.~Hutchinson,
``Deformation substructures and recrystallisation,''
\emph{Materials Science Forum}, vols.~558--559, pp.~13--22, 2007.

\bibitem{Ridha1982}
A.A.~Ridha and W.B.~Hutchinson,
``Recrystallisation mechanisms and the origin of cube texture in copper,''
\emph{Acta Metallurgica}, vol.~30, pp.~1929--1939, 1982.

\bibitem{Homma2003}
H.~Homma and B.~Hutchinson,
``Orientation dependence of secondary recrystallisation in silicon--iron,''
\emph{Acta Materialia}, vol.~51, pp.~3795--3805, 2003.

\bibitem{Lagneborg2014}
R.~Lagneborg, T.~Siwecki, S.~Zajac, and B.~Hutchinson,
\emph{The Role of Vanadium in Microalloyed Steels},
Swerea KIMAB Report KIMAB-2014-115, June 2014.
(Originally published in \emph{Scandinavian Journal of Metallurgy},
vol.~28, pp.~186--241, 1999; expanded edition available from Vanitec.)

\bibitem{Oscarsson1991}
A.~Oscarsson, W.B.~Hutchinson, and H.-E.~Ekstr\"om,
``Influence of initial microstructure on texture and earing in aluminium
sheet after cold rolling and annealing,''
\emph{Materials Science and Technology}, vol.~7, pp.~554--564, 1991.

\bibitem{Hutchinson2015}
B.~Hutchinson,
``Critical Assessment~16: Anisotropy in metals,''
\emph{Materials Science and Technology}, vol.~31, no.~12,
pp.~1393--1401, 2015.

\bibitem{Zajac1991}
S.~Zajac, T.~Siwecki, B.~Hutchinson, and M.~Attleg\aa{}rd,
``Recrystallization controlled rolling and accelerated cooling for high
strength and toughness in V--Ti--N steels,''
\emph{Metallurgical Transactions~A}, vol.~22A, pp.~2681--2694, 1991.

\bibitem{Hutchinson2011LUS}
B.~Hutchinson et~al.,
``Laser ultrasonics for process control in the metal industry,''
\emph{Nondestructive Testing and Evaluation}, vol.~26, nos.~3--4, 2011.

\bibitem{Hutchinson1999Phil}
B.~Hutchinson,
``Deformation microstructures and textures in steels,''
\emph{Philosophical Transactions of the Royal Society~A}, vol.~357,
no.~1756, pp.~1471--1485, 1999.

\bibitem{Hutchinson2025pssb}
B.~Hutchinson, M.~Malmstr\"om, A.~Jansson, and J.~L\"onnqvist,
``An experimental study of elasticity in polycrystalline iron,''
\emph{physica status solidi~(b)}, vol.~262, no.~6, 2025.

\bibitem{Hatherly1979}
M.~Hatherly and W.B.~Hutchinson,
\emph{An Introduction to Textures in Metals},
Institute of Metals, London, 1979; 2nd~ed.\ 1990.

\bibitem{Hutchinson2006Met}
B.~Hutchinson and J.~Hagstr\"om,
``Austenite decomposition structures in the Gibeon meteorite,''
\emph{Metallurgical and Materials Transactions~A}, vol.~37,
pp.~1811--1818, 2006.

\bibitem{Ray2011}
R.K.~Ray, B.~Hutchinson, and C.~Ghosh,
``Back-annealing of cold rolled steels through recovery and/or partial
recrystallisation,''
\emph{International Materials Reviews}, vol.~56, no.~2, pp.~73--97, 2011.

\bibitem{LagneborgZajac2001}
R.~Lagneborg and S.~Zajac,
``A model for interphase precipitation in V-microalloyed structural steels,''
\emph{Metallurgical and Materials Transactions~A}, vol.~32A, pp.~39--46,
2001.

\bibitem{Siwecki1995}
T.~Siwecki, B.~Hutchinson, and S.~Zajac,
``Recrystallisation controlled rolling of HSLA steels,''
Proc.\ Conf.\ Microalloying~'95, Pittsburgh, Iron \& Steel Soc., 1995,
pp.~197--211.

\bibitem{Siwecki1992}
T.~Siwecki and B.~Hutchinson,
``Modelling of microstructure evolution during recrystallisation
controlled rolling of HSLA steels,''
Proc.\ 33rd Mechanical Working and Steel Processing Conf., Warrendale,
ISS-AIME, 1992, pp.~397--406.

\bibitem{Hutchinson2014Bainite}
B.~Hutchinson, J.~Hagstr\"om, T.~Siwecki, J.~Komenda, R.~Lagneborg,
J.-E.~Hedin, and M.~Gladh,
``New vanadium-microalloyed bainitic 700~MPa strip steel product,''
\emph{Ironmaking and Steelmaking}, vol.~41, pp.~1--6, 2014.

\end{thebibliography}

\end{document}
